% =====================================================================================================
% Section VI: Conclusion. Writer: economics/conclusion writer (v2), 2026-09-24.
% Numbers and sources: paper/notes/writer2_economics.md (conclusion part).
% =====================================================================================================

\section{Conclusion}\label{sec:concl}

Optimization cuts both ways for measurement. Developers who choose model size and training data to minimize cost
generate data that say little about the technology they optimize over---along the expansion path, information about
curvature is fourth order in their allocation errors (Proposition~\ref{prop:ident})---and choices that say a good deal
about what they optimize for. Designed training experiments break the circle. IsoFLOP designs identify the elasticity of
substitution between parameters and data without a functional form (Proposition~\ref{prop:modelfree}), and the designs
of three laboratories agree on about 0.7. Given the technology, over-training reveals the value developers place on
compact models (Proposition~\ref{prop:wedge}). Under our reference technology, the median model in the clean sample of
open-weight releases is built as if that value, which for developers that serve their models is planned serving
expenditure, were three-quarters of lifetime cost; across open-weight releases, weighted by training compute, the share
rose from 0.04 in 2019--2022 to 0.82 in 2025.
\textcolor{red}{[TBD-m9: one sentence on what the controlled two-corpus experiment adds---$\sigma^*$ by corpus
(model-free and $\kappa$ free), the neutrality verdict for FineWeb-Edu versus FineWeb, and the direction of the high-$M$
extrapolation error.]}

\emph{What economists should use.}---For substitution between parameters and data in frontier-scale language models,
an elasticity of about 0.7: the model-free random-effects estimate across the large designs is 0.70, with a 95
percent confidence interval of 0.67 to 0.72, and parametric estimates that free the outer exponent agree with it there.
That the elasticity is below one is not a separate finding; it is what U-shaped IsoFLOP profiles imply
(Lemma~\ref{lemma:sigma}). Values of 0.5 to 0.6 from small-scale sweeps are heterogeneous and are best read as
design-specific. For mapping compute into loss, take the frontier elasticity $\gamma$ from the design whose compute range
and data are closest to the application; for frontier-scale language models the large designs give 0.14 to 0.18. The
allocation exponent and the compute-optimal ratio $M^*$ are not portable, and forecasts that depend on them, above all
of data demand, should carry their range across designs. The expenditure share $s$ is best used as an ordinal and
conditional measure. Its ranks and its trend are robust, and its sign is identified for 86 percent of the clean sample
under weak assumptions (Proposition~\ref{prop:pi}); its level depends on $M^*(C)$ at frontier scale, which no public
design observes, and it measures planned serving expenditure only for developers that serve their models.

\emph{Implications for scaling studies.}---Studies that aim to measure the technology rather than to tune a recipe
should place runs off the expansion path, estimate curvature from IsoFLOP profiles with the model-free estimator of
Proposition~\ref{prop:modelfree}, and test rather than impose $\kappa=1$. They should also report design-conditional
intervals and the uncertainty in $M^*(C)$, the object on which allocation decisions and economic forecasts depend most.

\emph{Limitations.}---Four caveats bound these conclusions. First, units: losses and token counts come from different
corpora, tokenizers, and counting conventions, so only unit-free objects compare across designs. Second, extrapolation:
88 percent of the clean sample use more compute than any design, and a quarter lie beyond the tokens per parameter
of the Farseer design. Where it can be checked, parametric extrapolation in tokens per parameter understates $w$, but nothing
is observed at frontier compute. \textcolor{red}{[TBD-m9: what the high-$M$ extrapolation test adds.]} Third,
disclosure: developers that do not report training tokens, including most closed developers, are absent, so the
revealed-demand results describe open-weight releases. Fourth, the economic applications inherit outside inputs: the
data-wall costs rest on repetition estimates from models of at most nine billion parameters, and loss is not
capability.

\emph{Extensions.}---The duality and identification arguments carry over to the isoquant between training and
test-time compute \citep{snell2024scaling,brown2024large}, whose elasticity growth models calibrate only loosely
\citep{erdil2025gate}, and to market structure, where a cost of quality that rises 50- to 160-fold per halving of
reducible loss is an endogenous sunk cost in the sense of \citet{sutton1991sunk}. More
broadly, designed training runs give production-function econometrics a rare setting in which what optimizing agents
reveal, and what they conceal, can be checked against experimental variation.
